\subsection{子向量丛、商向量丛、正合序列}

本号中，向量丛均指以 B 为基底的向量丛，向量丛态射均指 \(\mathrm{Id}_{\mathbf{B}}\)-态射。

7.5.1. 设 M 是一个向量丛。如果对每个点 \(b \in B\)，都存在 M 在 b 处的向量丛图 \(t = (\mathbf{U}, \varphi, \mathbf{E})\) 以及 E 的一个具有拓扑补空间的闭向量子空间 F，使得

\[
\varphi (\pi^ {- 1} (\mathrm{U}) \cap \mathrm{M} ^ {\prime}) = \mathrm{U} \times \mathrm{F}.
\]

在这些条件下，M′ 上存在唯一一个向量丛结构，使得从 M′ 到 M 的典范嵌入是一个态射。对每个 \(b \in B\)，M' 的纤维 \(M_{b}'\) 是闭向量子空间 \(M' \cap M_{b}\)，并且该空间包含于 \(M_{b}\)；M' 是 M 的闭子流形，并且 M' 的向量丛结构的底层流形结构与由 M 的流形结构诱导的结构一致。

7.5.2. 设 M' 是 M 的一个子向量丛。用 R 表示 M 的点 x、y 之间的如下关系：

存在 \(b\) 属于 \(B\)，使得 \(x \in M_b, y \in M_b\) 且 \(x - y \in M_b'\)。于是 \(R\) 是 \(M\) 上的正则等价关系（参见第 5.9.7 号）。集合 \(M/R\) 上存在唯一一个向量丛结构，使得从 \(M\) 到 \(M/R\) 的典范映射是一个态射。将其记为 \(M/M'\)，并称 \(M\) 对 \(M'\) 的商为如此定义的向量丛；对于 \(b\) 属于 \(B\) 的每个点，丛 \((M/M')_b\) 是商拓扑向量空间 \(M_b/M_b'\)，而 \(M/M'\) 上的流形结构是 \(M\) 上流形结构的商结构。

7.5.3. 保留 7.5.2 的假设。对于 B 的每个点 \(b_0\)，可以找到一个开邻域 U，其中包含 \(b_0\)；一个巴拿赫空间 F，它是两个闭子空间 \(F'\) 和 \(F''\) 的直和；以及一个向量丛同构 \(\iota\)，从 \(\mathbf{F}_{\mathbf{U}}\) 到 \(\mathbf{M}|\mathbf{U}\)，满足以下性质：

\begin{enumerate}
\def\labelenumi{(\roman{enumi})}
\item
  \(\iota'\) 是 \(\iota\) 在 \(\mathbf{U} \times \mathbf{F}'\) 上的限制，并且是从 \(\mathbf{F}_{\mathbf{U}}'\) 到 \(\mathbf{M}'|\mathbf{U}\) 的同构。
\item
  若 \(\rho\) 是从 M 到 \(\mathbf{M}'' = \mathbf{M} / \mathbf{M}'\) 的典范映射，且 \(\iota''\) 是 \(\iota\) 在 \(\mathbf{U} \times \mathbf{F}''\) 上的限制，则映射 \(\rho \circ \iota''\) 是从 \(\mathbf{F}_{\mathbf{U}}''\) 到 \(\mathbf{M}''|\mathbf{U}\) 的同构。
\end{enumerate}

7.5.4. 若向量丛态射 \(g: \mathbf{L} \to \mathbf{M}\) 的像包含于 \(\mathbf{M}'\)，则它是从 \(\mathbf{L}\) 到 \(\mathbf{M}'\) 的向量丛态射。

现在设 \(h: \mathbf{M} \to \mathbf{N}\) 是一个向量丛态射，并假定对 B 中的每个 \(b\)，\(h_b\) 在 \(\mathbf{M}_b'\) 上的限制为零。若 \(\rho\) 是从 M 到 \(\mathbf{M}/\mathbf{M}'\) 的典范态射，则存在唯一一个态射 \(\bar{h}\)，其从 \(\mathbf{M}/\mathbf{M}'\) 到 N，且满足 \(h = \bar{h} \circ \rho\)。

7.5.5. 设 P、Q 是两个向量丛，g 是从 P 到 Q 的态射；对 B 中的每个点 b，分别以 \(N_{b}\) 和 \(I_{b}\) 表示线性映射 \(g_{b}:P_{b}\to Q_{b}\) 的核和像。置 \(N=\bigcup_{b\in B}N_{b}\)，以及

\(I = \bigcup_{b \in B} I_b\)。称态射 \(g\) 是局部直的，当且仅当 \(N\) 是 P 的子向量丛且 I 是 Q 的子向量丛。于是态射 g 通过取商定义了从 P/N 到 I 的同构。称 N 为 g 的核，记为 Ker g。同样，子丛 I 称为 g 的像，记为 Im g。

若 \(r \geqslant \infty\)，则态射 \(g\) 局部直当且仅当 \(g\) 是子浸入。若 P 是有限秩的，则态射 \(g\) 局部直当且仅当 \(g\) 的向量秩局部常值，或者等价地，当且仅当 \(g\) 是子浸入。

7.5.6. 设 \(\mathbf{M} \xrightarrow{f} \mathbf{M}' \xrightarrow{g} \mathbf{M}''\) 是两个向量丛态射。称序列 \((f, g)\) 是局部直正合的，当且仅当 \(f\) 和 \(g\) 都局部直且 \(\operatorname{Im} f = \operatorname{Ker} g\)。若 \(g \circ f = 0\)，令 \(\mathbf{D}\) 为满足以下条件的点 \(b \in \mathbf{B}\) 的集合：序列 \(\mathbf{M}_b \xrightarrow{f_b} \mathbf{M}_b' \xrightarrow{g_b} \mathbf{M}_b''\) 是直正合的（即 \(\operatorname{Ker} f_b\) 和 \(\operatorname{Im} g_b\) 有拓扑补空间，且 \(\operatorname{Im} f_b = \operatorname{Ker} g_b\)）。则 D 是开集，并且序列 \(\mathbf{M}|\mathbf{D} \xrightarrow{f} \mathbf{M}'|\mathbf{D} \xrightarrow{g} \mathbf{M}''|\mathbf{D}\) 是局部直正合的。

任意长度的局部直正合序列也以同样方式定义。语言上滥用时，有时用“直正合序列”代替“局部直正合序列”。

7.5.7. 设 \(0 \to \mathbf{M} \xrightarrow{f} \mathbf{M}' \xrightarrow{g} \mathbf{M}'' \to 0\) 是一个向量丛态射序列。这个序列局部直正合的充要条件是：\(f\) 是从 \(\mathbf{M}\) 到子向量丛 \(f(\mathbf{M})\) 的同构，该子丛属于 \(\mathbf{M}'\)，并且 \(g\) 通过取商定义了从商向量丛 \(\mathbf{M}'/f(\mathbf{M})\) 到 \(\mathbf{M}''\) 的同构。

